\section{Introduction}\label{sec:introduction}

The entries of a matrix's first column are a row apart on the input tape
and consecutive in the output. A transposer may store coded information
as it works, and its head movements may depend on what it has read. If
we fix those movements, we can trace the possible dependencies of an
output bit. Fixing the movements also restricts the input. The question
is how much information survives the restriction.

For a fixed tape trajectory, a depth-reduction theorem supplies a small
set of steps whose values can be revealed so that the remaining
dependency paths have bounded length \cite[Corollary~O]{PaperFourteen}.
On the inputs sharing both the trajectory and the revealed values, each
output can then be recovered from a bounded set of input coordinates.
The recovered coordinate and its input copy can be far apart in the
incidence graph of that relation. Such separation is compatible with
binary storage codes of rate above one half
\cite[Corollary~A$'$]{PaperThirteen}. A linear-time transposer would
require rates arbitrarily close to one.

\Cref{thm:main} gives the degree dependence of this requirement. For
every sufficiently large degree cap it produces a family with rates
bounded below while its minimum paired distance tends to infinity. The
coordinates are bits of packets. A coding upper bound with a larger
deficit would rule out the assumed machine. \Cref{cor:gap} records that
conditional test, including success probability when the machine is
correct only on some inputs. The construction below shows why the test
cannot succeed for unrestricted storage codes.

\paragraph{The abstract coding requirement is attainable.}
Starting from the above-half-rate linear family of
\cite[Corollary~U(a,b)]{PaperThirteen}, we construct, at every sufficiently
large cap $\Delta$, admissible binary families of rate at least
$1-C\Delta^{-\alpha}$, with fixed $C>0$ and $0<\alpha<1/2$.
\Cref{thm:majority} amplifies a fixed subspace by a signed sum of tensor
operators. Odd parity preserves the absence of short odd closed walks.
Taking a syndrome fibre inside the binary cube converts the resulting
rank bound into a binary rate bound without adding coordinates.
\Cref{thm:construction} gives the constants and the order of quantifiers.

The same encoder partitions the entire cube before a source subset is
chosen. \Cref{thm:all-length} extends these partitions to every
sufficiently large length, retaining every coordinate and bounding the
input mass in small fibres. These properties therefore cannot, by
themselves, distinguish a tape-induced code from an abstract code.
\Cref{cor:abstract-gap} rules out the unrestricted slow-gap premises of
\cref{cor:gap}. A lower bound must use an additional restriction imposed
by a physical computation.

A second question concerns the individual bit positions of a physical
machine. \Cref{thm:simulator} groups source transitions into batches but
executes each head move and each append on an ordinary bit machine.
\Cref{thm:native-tail} uses its actual append times. The packet graph connects every bit of a permitted helper packet to
every bit of its target packet. The native graph follows each physical
append backwards through the simulator. On a set $G$ of correctly
transposed inputs, the native theorem bounds the input mass in fibres
whose rate falls appreciably below $\log_2|G|/L$. Each fibre keeps nearly
all of its trace's original bit positions. Revealing complete simulator batches and counting their joint record
support gives the same rate-loss exponent as the packet theorem.

Finally, \cref{sec:patches} asks how much a local set of helpers can tell
us about the bit it is meant to recover. The answer accounts separately
for a bit's own bias and the information supplied by its helpers. A
fixed advantage over chance on a positive fraction of separated targets
requires increasingly large patch overlap as the code rate approaches
one. The outstanding question is whether a linear-time tape computation
would have to provide such an advantage at smaller overlap. The
construction supplies no tape schedule or efficient uniform tape
algorithm for generating coefficients and evaluating its encoder and
recoveries.

\paragraph{Relation to earlier work.}
The finite-control simulation of Hartmanis and Stearns evaluates
consecutive tape operations from a bounded neighbourhood
\cite[Theorem~2]{HartmanisStearns}. Hopcroft, Paul and Valiant synchronize
tape-block crossings \cite[\S2]{HPV1977}. Williams describes the
computation graph of an aligned machine using a constant number of
movement bits per block \cite[\S\S2--3]{Williams}. Here the sequence of batch descriptors determines the in-batch append
flags and when the terminal drain starts. The common output length
then fixes the number of bits appended during the drain. Earlier in this series, an external-memory
simulation fixed an event word at the scale of entire rows
\cite[Theorem~EM]{PaperEleven}. Our batch size is fixed before the
matrix dimension grows. Conditional entropy as a measure of imperfect
local recovery appears in earlier work on recoverable systems
\cite[\S IV]{ElishcoBarg}. Our prediction bound uses the separated-port
argument of \cite[Lemma~4.1]{PaperFive} and the soft-recovery inequality
for nested elimination \cite[Theorem~S]{SoftRecovery}.
The construction uses the odd-field linear code and the symmetrized
graphs of \cite[Corollary~U(a,b)]{PaperThirteen}. Its binary conversion
uses fibres of a linear syndrome restricted to binary inputs. The
tensor majority argument uses only a basis of a fixed subspace and an
arbitrary complementary basis; no diagonalization is assumed.
Binary storage rates approaching one on triangle-free graphs were
already obtained by Barg, Schwartz and Yohananov
\cite[Theorem~2]{BSY2024}. Golovnev and Haviv construct low-minrank
graphs excluding all odd cycles up to any prescribed fixed length
\cite[Theorem~22]{GolovnevHaviv2021}. Here the degree cap is fixed
before the coordinate count and minimum odd girth grow; only then
do we compare different caps. The graph containing the amplified
support is an instance of the standard non-complete extended
$p$-sum (NEPS) product, with all odd-weight coordinate-change masks
\cite[\S2.3 and Theorem~3.1]{NEPS2011}. Our fitting matrix need not
be that graph's adjacency matrix.
\Cref{sec:barriers} explains why literal tensor checks and a specified
two-slot local-code construction do not supply the same amplification.

\paragraph{Formal proofs.}
The Lean declarations named with each formalized result are its proofs.
The English outlines explain their constructions, estimates and parameter
choices. The amplification theorems have an explicit seed hypothesis;
\cref{prop:seed-exists} supplies it in a written proof from the earlier
paper. That earlier construction's arithmetic-geometric existence
arguments are not formalized here. \Cref{sec:verification} records this
boundary and the verification procedure. The declaration index identifies
the source modules. We first give the tape requirement, then the abstract
construction and its barriers, and finally the simulator, native
partitions and physical prediction bounds.
